% !TEX program = lualatex
\documentclass[aspectratio=43,professionalfonts,handout]{beamer}
\usepackage{beamer_template}

\newcommand{\LectureNum}{27}
\newcommand{\LectureTitle}{加法定理の応用（総合演習）}
\newcommand{\TextPages}{p.\,168-170（練習問題3-A, 3-B）}
\newcommand{\NextTopic}{後期前半の総整理}
\newcommand{\NextPages}{p.\,102-139}
\newcommand{\CourseName}{基礎数学B}
\renewcommand{\TermName}{後期}

\begin{document}

% -----------------------------------------------
% スライド1: 表紙＋目標
% -----------------------------------------------
\begin{frame}{\CourseName\quad 第\LectureNum 回「\LectureTitle 」}
  {\small \TermName\quad \TeacherName\quad ／\quad 教科書 \TextPages\quad
  \textbf{目標}：加法定理・2倍角・積和・合成を組み合わせて問題を解ける}

\end{frame}

% -----------------------------------------------
% スライド2: 公式・性質
% -----------------------------------------------
\begin{frame}{公式・性質}
  \begin{block}{}
  \begin{itemize}
    \item 加法定理・2倍角・半角・合成の公式一覧
  \end{itemize}
  \end{block}
\end{frame}

% -----------------------------------------------
% まとめ
% -----------------------------------------------
\begin{frame}{まとめ}
  \begin{itemize}
    \item 加法定理・2倍角・半角・合成を使い分けて問題を解く
    \item 式変形では次数を下げる方向（半角公式，合成）を意識する
    \item 定義域を確認して解の範囲を正しく求める
  \end{itemize}
\end{frame}

\end{document}
